\begin{textsample}{POS Dim 2 – 2023\_c5 – Score -1.00 – 2023\_c5/enade\_2023\_med\_c56e7c1ea6.txt}  \label{ex:f2_pos_262}
CAPÍTULO 4 DISTRIBUIÇÃO DE CONCEITOS A finalidade deste capítulo é apresentar um panorama nacional sobre a distribuição do Conceito Enade 2023.
Para traçar esse panorama, recorre -se a tabelas e gráficos, que são acompanhados de comentários e análises que contribuem para a contextualização dos dados neles expostos.
Para visualizar o panorama do Enade 2023, em nível nacional, na primeira seção deste capítulo, os dados sobre o número de cursos participantes são organizados por grande região.
Na segunda, o mapeamento do conceito do exame é feito por categoria administrativa das IES e modalidade de oferta dos cursos, de maneira cruzada às grandes regiões.
Na última seção, a lógica de apresentação dos dados é feita por organização acadêmica e, também, por grande região de funcionamento dos cursos, no país como um todo. 4.1 PANORAMA NACIONAL DA DISTRIBUIÇÃO DOS CONCEITOS A lógica adotada na apresentação dos dados, nesta seção e nas duas seguintes, pressupõe o registro de duas informações de natureza reiterativa.
O primeiro diz respeito aos valores possíveis do Enade, que variam em níveis de 1 a 5, sendo que, à medida que esses valores aumentam, o curso é melhor avaliado, e ao fato de que aqueles cursos que tiveram somente um ou nenhum concluinte participante não obtêm o Conceito Enade, ficando sem conceito ( SC ), como ressaltado na apresentação deste relatório.
O segundo registro é que, caso haja diferença entre os cursos tabulados neste capítulo e no capítulo 1, ela corresponderá aos cursos que não foram avaliados, ou seja, àqueles, em princípio, sem estudantes concluintes que tivessem sido inscritos para a prova.
Na Tabela 4.1, são apresentadas a quantidade e a distribuição de cursos de Medicina participantes do Enade 2023, por faixa de conceito e grande região.
A partir dos dados da Tabela 4.1, nota -se que, dos 309 cursos participantes, 104 ( 33,7 \% ) classificaram -se com conceito 3, o valor modal.
Esse foi também o conceito modal na maioria das regiões, exceto a região Sul : Centro-Oeste ( 40,0 \% ), Nordeste ( 33,3 \% ), Norte ( 54,2 \% ) e Sudeste ( 30,2 \% ).
O conceito 4 foi o segundo mais frequente em nível nacional ( 30,4 \%, correspondendo a 94 cursos ), tendo sido o segundo conceito mais frequente em todas as regiões.
O conceito 2 foi o terceiro ( 15,2 \%, correspondendo a 47 cursos ).
Houve, ainda, 45 ( 14,6 \% ) cursos que receberam o conceito 5 e 15 ( 4,9 \% ) cursos que receberam o conceito 1.
Dos 309 cursos de Medicina, 4 ( 1,3 \% ) ficaram sem conceito ( SC ).
A região Norte participou com 24 cursos, o que corresponde a 7,8 \% do total nacional.
Desses, 13 cursos receberam o conceito 3, o que equivale a 54,2 \% do total regional, o conceito modal na região, conforme mencionado.
A cinco cursos ( 20,8 \% ) atribuiu -se o conceito 4 ; a três cursos ( 12,5 \% ), o conceito 2 ; e a três cursos ( 12,5 \% ) o conceito 1.
Ainda nessa região, nenhum curso ( 0,0 \% ) obteve conceito 5 e nenhum curso ( 0,0 \% ) ficou sem conceito ( SC ).
A região Nordeste participou com 75 cursos, correspondentes a 24,3 \% do total nacional.
Nessa região, 25 cursos ( 33,3 \% em termos regionais ) obtiveram conceito 3, o conceito modal para a região, como já comentado.
Os conceitos 4 e 2 foram atribuídos à 23 e a 13 cursos ( 30,7 \% e 17,3 \%, respectivamente ).
Ainda nessa região, 12 cursos ( 16,0 \% ) foram avaliados com o conceito 5 e dois cursos ( 2,7 \% ) foram avaliados com o conceito 1.
Nenhum curso ( 0,0 \% ) ficou sem conceito ( SC ).
Para os 126 cursos participantes da região Sudeste ( 40,8 \% do total nacional ), o conceito 3 foi o conceito modal, como já anteriormente assinalado, tendo sido obtido por 38 cursos ( 30,2 \% ).
O conceito 4 foi atribuído a 34 cursos ( 27,0 \% ) e o conceito 2, a 23 cursos ( 18,3 \% ).
Receberam os conceitos 5 e 1, respectivamente, 20 cursos ( 15,9 \% ) e nove cursos ( 7,1 \% ).
Nessa região, dois cursos ( 1,6 \% ) ficaram sem conceito ( SC ).
Os 54 cursos da região Sul corresponderam a 17,5 \% do total nacional.
Como já indicado, houve predominância do conceito 4 ( conceito modal ), atribuído a 44,4 \% dos cursos da região, ou seja, atribuído a 24 dos 54 cursos participantes na região Sul.
O conceito 3 foi atribuído a 16 cursos ( 29,6 \% ) ; o conceito 5, oito cursos ( 14,8 \% ) ; e o conceito 2, a quatro cursos ( 7,4 \% ).
Na região Sul, nenhum curso ( 0,0 \% ) obteve conceito 1 e dois cursos ( 3,7 \% ) ficaram sem conceito ( SC ).
Já dos 30 cursos participantes na região Centro-Oeste ( 9,7 \% do total nacional ), 12 ( 40,0 \% em termos regionais ) receberam conceito 3, o conceito modal.
Oito cursos ( 26,7 \% ) obtiveram o conceito 4, cinco cursos ( 16,7 \% ) obtiveram o conceito 5 e quatro cursos ( 13,3 \% ), o conceito 2.
Nessa região, um curso ( 3,3 \% ) obteve o conceito 1 e nenhum curso ( 0,0 \% ) ficou sem conceito ( SC ).
Em um gráfico de distribuição cumulativa, em cada valor de abscissa ( eixo x ), o valor apresentado no eixo das ordenadas ( eixo y ) é o valor acumulado da variável ( Conceito Enade ).
Em linhas gerais, em um gráfico desta mesma tipologia, poligonais mais à esquerda representam, no caso dos dados apresentados nesta seção, as grandes regiões com a pior distribuição, e poligonais mais à direita, as grandes regiões com os melhores conceitos.
Respectivamente, 20 cursos ( 15,9 \% ) e nove cursos ( 7,1 \% ).
Nessa região, dois cursos ( 1,6 \% ) ficaram sem conceito ( SC ).
Os 54 cursos da região Sul corresponderam a 17,5 \% do total nacional.
Como já indicado, houve predominância do conceito 4 ( conceito modal ), atribuído a 44,4 \% dos cursos da região, ou seja, atribuído a 24 dos 54 cursos participantes na região Sul.
O conceito 3 foi atribuído a 16 cursos ( 29,6 \% ) ; o conceito 5, oito cursos ( 14,8 \% ) ; e o conceito 2, a quatro cursos ( 7,4 \% ).
Na região Sul, nenhum curso ( 0,0 \% ) obteve conceito 1 e dois cursos ( 3,7 \% ) ficaram sem conceito ( SC ).
Já dos 30 cursos participantes na região Centro-Oeste ( 9,7 \% do total nacional ), 12 ( 40,0 \% em termos regionais ) receberam conceito 3, o conceito modal.
Oito cursos ( 26,7 \% ) obtiveram o conceito 4, cinco cursos ( 16,7 \% ) obtiveram o conceito 5 e quatro cursos ( 13,3 \% ), o conceito 2.
Nessa região, um curso ( 3,3 \% ) obteve o conceito 1 e nenhum curso ( 0,0 \% ) ficou sem conceito ( SC ).
Em um gráfico de distribuição cumulativa, em cada valor de abscissa ( eixo x ), o valor apresentado no eixo das ordenadas ( eixo y ) é o valor acumulado da variável ( Conceito Enade ).
Em linhas gerais, em um gráfico desta mesma tipologia, poligonais mais à esquerda representam, no caso dos dados apresentados nesta seção, as grandes regiões com a pior distribuição, e poligonais mais à direita, as grandes regiões com os melhores conceitos.
A partir destas diretrizes, no Gráfico 4.1, é apresentada a distribuição cumulativa do Conceito Enade, segundo a grande região.
No Gráfico 4.1, observando a linha vertical pontilhada, que ajuda a visualizar o percentual de conceitos acumulados até o valor 3, ou seja, conceitos 1, 2 e 3, pode -se notar, por exemplo, que a região Norte ( linha azul ) apresenta 79,2 \% dos cursos com conceito 3 ou abaixo deste valor.
A região Sul ( linha roxa ), por outro lado, apresenta 38,5 \% dos cursos na mesma situação : com conceito 3 ou abaixo.
Destaca -se que, a região Norte ( linha azul ), o conceito 5 não foi alcançado por nenhum curso, logo os cursos têm conceito menor ou igual a 4, e a respectiva poligonal cumulativa alcança 100 \% no conceito 4.
A região Norte apresenta os cursos com conceitos mais baixos, ao passo que a região Sul apresenta a distribuição com valores maiores.
Os cursos em IES situadas na região Sudeste ( linha vermelha ) apresentam uma situação intermediária. 4.2 CONCEITOS POR CATEGORIA ADMINISTRATIVA, MODALIDADE DE OFERTA DOS CURSOS E GRANDE REGIÃO Na Tabela 4.2, é apresentada a distribuição dos cursos de Medicina participantes do Enade 2023, por categoria administrativa e modalidade de oferta, de acordo com os conceitos por eles alcançados, segundo as grandes regiões brasileiras.
De acordo com as informações da Tabela 4.2, dos 309 cursos participantes, 121 ( 39,2 \% ) eram ministrados em instituições públicas e 188 ( 60,8 \% ), em instituições privadas.
Quanto à modalidade de oferta, todos os 309 eram presenciais.
De acordo com as informações da mesma tabela, em termos nacionais, entre instituições públicas e privadas, observa -se que, dos 45 cursos avaliados com conceito 5, 37 foram oferecidos em IES públicas, e 8, em IES privadas.
Dos 188 cursos participantes de IES privadas, o conceito 3 foi o valor modal, atribuído a 82 ( 43,6 \% ) cursos.
Em relação aos demais cursos de instituições privadas, 46 cursos ( 24,5 \% ) foram avaliados com o conceito 4, 41 cursos ( 21,8 \% ) foram avaliados com o conceito 2, dez cursos ( 5,3 \% ) foram avaliados com o conceito 1 e, oito cursos ( 4,3 \% ), com o conceito 5.
Nessa categoria, 1 curso ( 0,5 \% ) ficou sem conceito ( SC ).
Na rede pública, o conceito modal foi o 4, atribuído a 48 cursos, correspondentes a 39,7 \% dos 121 cursos da categoria.
Dos demais cursos participantes, 37 cursos ( 30,6 \% ) obteve o conceito 5, 22 cursos ( 18,2 \% ) receberam o conceito 3 e seis cursos ( 5,0 \% ), o conceito 2.
O conceito 1 foi atribuído a cinco cursos ( 4,1 \% ).
Nessa categoria administrativa, três cursos ( 2,5 \% ) ficaram sem conceito ( SC ).
Dos 309 cursos participantes, o conceito 3 foi o valor modal, atribuído a 104 ( 33,7 \% ) cursos.
Dos demais cursos, 94 cursos ( 30,4 \% ) receberam o conceito 4 ; 47 cursos ( 15,2 \% ) receberam o conceito 2 ; 45 cursos ( 14,6 \% ), o conceito 5 ; 15 cursos ( 4,9 \% ), o conceito 1 ; e quatro cursos ( 1,3 \% ) ficaram sem conceito ( SC ).
Na análise por região, observa -se, que, na região Norte, as instituições públicas participaram com 12 cursos ( 50,0 \% total regional ), dos quais nenhum ficou sem conceito ( SC ) ou recebeu o conceito 5.
O conceito modal foi 3, atribuído a cinco cursos, correspondendo a 41,7 \% dessa categoria na região.
O conceito 4 teve uma concentração muito próxima do modal, com quatro curso ( 33,3 \% ), e o conceito 2 teve três curso ( 25,0 \% ).
Os conceitos 5 e 1 não foram atribuídos a nenhum curso.
As instituições privadas participaram com 12 cursos na região Norte ( 50,0 \% do total regional ), e nenhum curso ficou sem conceito ( SC ).
O conceito modal nas instituições privadas na região foi o 3, atribuído a oito cursos, o que corresponde a 66,7 \% dessa categoria na região.
Os conceitos 1, 4 e 5 foram atribuídos a três, um e nenhum curso, respectivamente.
O conceito 2 foi atribuído a nenhum curso ( 0,0 \% ).
Nenhum curso ficou sem conceito ( SC ) ou recebeu os demais conceitos.
Dos cursos oferecidos na região Norte, nenhum ficou sem conceito ( SC ).
Dos 24 cursos oferecidos, 13 obtiveram o conceito 3 ( modal ) ; cinco, o conceito 4 ; três, o conceito 2 ; três, o conceito 1 ; e nenhum, o conceito 5.
Na região Nordeste, a rede privada concentra 37 dos 75 cursos participantes, o equivalente a 49,3 \% do total da região.
Entre os cursos oferecidos por IES privadas no Nordeste, 19 cursos foram avaliados com conceito 3, o conceito modal. 12 cursos obteve o conceito 2 ; cinco cursos, o conceito 4 ; um curso, o conceito 1 ; e nenhum curso obteve o conceito 5.
Nessa categoria administrativa e grande região, nenhum curso ficou sem conceito ( SC ).
As instituições públicas da região Nordeste participaram com 38 cursos ( 50,7 \% do total da região ).
Desses, 18 ( 47,4 \% ) obtiveram conceito 4, o valor modal, seguidos por 12 cursos, que receberam o conceito 5.
O conceito 3 foi atribuído a seis cursos ; o conceito 2, a um curso ; e o conceito 1, a um curso.
Nenhum curso ficou sem conceito ( SC ).
Dos cursos oferecidos na região Nordeste : nenhum ficou sem conceito ( SC ), dois obtiveram o conceito 1 ; 13, o conceito 2 ; 25, o conceito 3, o valor modal ; 23, o conceito 4 ; e 12, o conceito 5.
Na região Sudeste, a proporção de cursos da rede privada, 73,8 \%, correspondeu a 93 dos 126 cursos participantes.
Nessa combinação de categoria administrativa e grande região, o conceito modal foi o 3, atribuído a 35 cursos ; e nenhum curso ficou sem conceito ( SC ).
Os demais foram avaliados com os seguintes conceitos : 4 ( 24 cursos ), 2 ( 23 cursos ), 1 ( seis cursos ) e 5 ( cinco cursos ).
Dos 33 cursos oferecidos por instituições públicas na região Sudeste, 15 obtiveram o conceito 5 ; dez, o conceito 4 ; três, o conceito 3 ; nenhum obteve conceito 2 e três obtiveram o conceito 1.
Dois cursos ficaram sem conceito ( SC ).
As instituições públicas na região Sudeste participaram com 33 cursos, aos quais foram atribuídos os conceitos 1 ( três cursos ), 2 ( nenhum curso ), 3 ( três cursos ), 4 ( dez cursos ) e 5 ( quinze cursos ).
Dois cursos da região Sudeste oferecido por IES públicas ficaram sem conceito ( SC ).
As instituições privadas concentraram 35 dos 54 cursos participantes da região Sul, 64,8 \% do total regional.
Desses, 14 cursos obtiveram conceito 3, o conceito modal.
Nessa combinação de categoria administrativa e grande região, 13 receberam o conceito 4 ; quatro receberam o conceito 2 ; três, o conceito 5 ; e nenhum obteve o conceito 1.
Um curso ficou sem conceito ( SC ).
As instituições públicas na região Sul participaram com 19 cursos ( 35,2 \% dos cursos da região ), aos quais foram atribuídos os conceitos 1 ( nenhum curso ), 2 ( nenhum curso ), 3 ( dois cursos ), 4 ( 11 cursos ) e 5 ( cinco cursos ).
Um curso da região Sul oferecido por IES públicas ficou sem conceito ( SC ).
Na região Centro-Oeste, 11 dos 30 cursos participantes eram de IES privadas ( 36,7 \% em termos regionais ).
Desses, seis cursos receberam conceito 3, o valor modal.
Dos demais cursos, três receberam o conceito 4, dois receberam o conceito 2, nenhum recebeu o conceito 5 e nenhum recebeu o conceito 1.
Além disso, nenhum curso ficou sem conceito ( SC ).
Quanto aos 19 cursos oferecidos por instituições públicas na região Centro-Oeste ( 63,3 \% do total regional ), um ficou com conceito 1 ; dois, com conceito 2 ; seis, com conceito 3 ; cinco, com conceito 4 ; e cinco ficaram com conceito 5.
Nenhum ficou sem conceito ( SC ) nessa região.
Nesse gráfico, poligonais mais à esquerda representam Conceito Enade mais baixos, e poligonais mais à direita, Conceito Enade mais altos.
É importante observar que, para cada poligonal presente nesse gráfico, a porcentagem constante no eixo y refere -se ao universo de cursos oferecidos por instituições públicas ( poligonal azul ) ou privadas ( poligonal verde ), respectivamente.
Os cursos em IES públicas ( poligonal azul ) apresentam uma poligonal mais à direita do que os em IES privadas ( poligonal verde ) e, consequentemente, uma distribuição de Conceito Enade com valores maiores. 4.3 CONCEITOS POR ORGANIZAÇÃO ACADÊMICA E GRANDE REGIÃO Na Tabela 4.3, encontra -se a distribuição dos conceitos atribuídos aos cursos participantes do Enade 2023, na área de Medicina, por organização acadêmica, segundo as grandes regiões brasileiras.
Como mostram os dados da Tabela 4.3, dos 309 cursos de Medicina participantes, 187 foram oferecidos em universidades, 67 em centros universitários e 55 em faculdades ; nenhum curso foi oferecido por CEFET/IF.
Essa distribuição corresponde a, respectivamente, 60,5 \%, 21,7 \%, 17,8 \% e 0,0 \% dos cursos.
De acordo com os dados apresentados na mesma tabela, de todos os 104 cursos avaliados com o conceito 3, 50 eram vinculados a universidades.
Os cursos vinculados a esse tipo de organização acadêmica tiveram o conceito 4 como conceito modal, atribuído a 71 cursos ( 38,0 \% ).
Os demais cursos ofertados por universidades receberam os conceitos 5 ( 33 cursos ), 2 ( 23 cursos ) e 1 ( sete cursos ).
Três cursos ficaram sem conceito ( SC ).
Dos cursos vinculados a centros universitários, o conceito modal também foi o 3, atribuído a 33 cursos ( 49,3 \% ).
Um curso vinculado a esse tipo de organização acadêmica ficou sem conceito ( SC ), e os outros cursos receberam os conceitos 4 ( 14 cursos ), 2 ( 11 cursos ), 5 ( quatro cursos ) e 1 ( quatro cursos ).
Dos 55 cursos mantidos por faculdades, 21 ( 38,2 \% ) receberam conceito 3, o conceito modal.
O segundo conceito mais frequentemente atribuído aos cursos oferecidos por essa categoria de IES foi o conceito 2, recebido por 13 cursos.
Na sequência, nove cursos obtiveram conceito 4 ; oito cursos, o conceito 5 ; e quatro cursos, o conceito 1.
Dos cursos oferecidos em faculdades, nenhum curso ficou sem conceito ( SC ).
Considerando -se, separadamente, as regiões brasileiras, verifica -se que, na região Norte, as universidades participaram com 13 cursos avaliados, aos quais foram atribuídos os seguintes conceitos : 1 a um curso ; 2 a três cursos ; 3, o conceito modal, a cinco cursos ; conceito 4 a quatro cursos ; e 5 a nenhum curso.
Nenhum curso ficou sem conceito ( SC ).
Os centros universitários da região Norte foram representados por oito cursos, dos quais dois cursos obteve conceito 1 ; nenhum curso, conceito 2 ; cinco cursos, conceito 3, o valor modal ; um curso, o conceito 4 ; e nenhum curso, o conceito 5.
Nenhum curso ficou sem conceito ( SC ).
As faculdades participaram com três cursos na região Norte, os quais receberam o conceito 3 Na região Nordeste, as universidades participaram com 45 dos 75 cursos da área de Medicina oferecidos.
O conceito modal foi 4, atribuído a 20 cursos.
Outros 12 cursos obtiveram o conceito 5 ; oito cursos, o conceito 3 ; quatro cursos, o conceito 2 ; e um curso, o conceito 1.
Nenhum curso ficou sem conceito ( SC ).
Os centros universitários contaram com 15 cursos participantes na região Nordeste, dez dos quais receberam o conceito modal 3.
Para os demais cursos, foram atribuídos os conceitos 2 ( quatro cursos ) e 4 ( um curso ).
As faculdades foram representadas por 15 cursos na região Nordeste, os quais receberam os conceitos 1 ( um curso ), 2 ( cinco cursos ), 3 ( sete cursos ), 4 ( dois cursos ) e 5 ( nenhum curso ).
Nenhum curso Ficou sem conceito ( SC ).
Na região Sudeste, as universidades concentraram 65 dos 126 cursos de Medicina da região.
Dos cursos oferecidos em IES com este tipo de organização acadêmica, nessa região, o conceito modal foi 4, atribuído a 20 cursos.
Os demais cursos receberam os conceitos 3 ( 17 cursos ), 2 ( 12 cursos ), 5 ( dez cursos ) e 1 ( quatro cursos ).
Dois cursos ficaram sem conceito ( SC ).
Os centros universitários participaram com 31 cursos na região Sudeste, dos quais 12 obtiveram o conceito modal 3.
O conceito 4 foi atribuído a nove cursos ; o conceito 2 a cinco cursos ; o conceito 5 a três cursos ; e o conceito 1 a dois cursos.
Nenhum curso ficou sem conceito ( SC ) nessa combinação de organização acadêmica e região.
Já as faculdades foram representadas por 30 cursos na região Sudeste, dos quais nenhum ficou sem conceito ( SC ).
O conceito modal 3 foi atribuído a 9 cursos.
Aos demais cursos foram atribuídos os conceitos 5 ( sete cursos ), 2 ( seis cursos ), 4 ( cinco cursos ) e 1 ( três cursos ).
Dos 54 cursos da região Sul, 44 eram mantidos por universidades, para os quais o conceito modal foi 4, atribuído a 20 cursos.
Os demais cursos receberam os conceitos 3 ( 12 cursos ), 5 ( sete cursos ), 2 ( quatro cursos ) e 1 ( nenhum curso ).
Em IES com esse tipo de organização acadêmica, um curso da região Sul ficou sem conceito ( SC ).
Dos seis cursos participantes de centros universitários da região Sul, um ficou sem conceito ( SC ), e dois cursos receberam conceito 4, o valor modal.
Os conceitos restantes ficaram assim distribuídos : 3 ( dois cursos ), 5 ( um curso ), 2 ( nenhum curso ) e 1 ( nenhum curso ).
Quanto aos quatro cursos vinculados a faculdades na região Sul, dois cursos receberam o conceito 4, valor modal.
Os demais conceitos ficaram assim distribuídos : 3 ( dois cursos ), 5 ( nenhum curso ), 2 ( nenhum curso ) e 1 ( nenhum curso ).
Dos quatro cursos, nenhum ficou sem conceito ( SC ).
Na região Centro-Oeste, 20 dos 30 cursos oferecidos eram mantidos por universidades.
Para esse tipo de organização acadêmica, o conceito modal foi 3, atribuído a oito cursos.
Em relação às universidades, nenhum curso da região ficou sem conceito ( SC ).
Os demais cursos receberam os conceitos 4 ( sete cursos ), 5 ( quatro cursos ), 1 ( um curso ) e 2 ( nenhum curso ).
Os centros universitários da região Centro-Oeste contaram com sete cursos : quatro cursos foram avaliados com o conceito 3 ( modal ) ; dois cursos, com o conceito 2 e um curso, com o conceito 4.
Dos três cursos oferecidos em faculdades na região Centro-Oeste, dois cursos receberam o conceito 2 ( modal ) e um curso, o conceito 5.
De maneira semelhante à forma como se procedeu em relação à seção anterior, esta seção encerra -se com o Gráfico 4.3, que apresenta a distribuição do Conceito Enade dos cursos de Medicina segundo a organização acadêmica da IES.
No Gráfico 4.3, apresenta -se a distribuição do Conceito Enade segundo a organização acadêmica da IES.
A leitura desse último gráfico é semelhante à dos anteriores, seguindo -se a lógica de que poligonais deslocadas para a esquerda correspondem aos conceitos mais baixos ( 1, 2 e 3 ) e poligonais deslocadas para a direita correspondem a Conceitos Enade mais altos.
Nas universidades ( linha azul ), todos os cursos obtiveram conceito 5 ou abaixo disso, já que a poligonal cumulativa alcança 100 \% para o conceito 5.
Além disso, os cursos de universidades apresentam uma poligonal mais à direita do que as demais, denotando uma melhor distribuição de Conceitos Enade.
No outro extremo, encontram -se os cursos mantidos por faculdades ( linha vermelha ), apresentando as respectivas poligonais mais à esquerda e, consequentemente, uma distribuição de Conceitos Enade com os menores valores.
Os cursos mantidos por centros universitários ( linha verde ) apresentam poligonal mais central e, consequentemente, uma distribuição de Conceitos Enade com os valores intermediários.

% matched lemmas: 
\end{textsample}
